%\documentclass[11pt,a4paper,uplatex,dvipdfmx]{ujarticle} 		% for uplatex
\documentclass[11pt,a4j,dvipdfmx]{jarticle} 					% for platex
\input{pieces/form00_header} % pieces
\input{pieces/kakenhi7} % pieces
\input{pieces/form01_dcpd_header} % pieces
\input{pieces/hook3} % pieces
%#Name: cpd
\input{pieces/form03_dcpd_headers} % pieces
\input{pieces/form04_cpd_header} % pieces
% ===== Global definitions for the Kakenhi form ======================
%　基本情報
%
%------ 研究課題名  -------------------------------------------
\newcommand{\研究課題名}{象の卵}

%----- 研究機関名と研究代表者の氏名-----------------------
\newcommand{\研究機関名}{逢坂大学}
\newcommand{\研究代表者氏名}{湯川秀樹}
\newcommand{\me}{\underline{\underline{H.~Yukawa}}} 
\input{pieces/inst_dcpd} % pieces
% user07_header
% ===== my favorite packages ====================================
% ここに、自分の使いたいパッケージを宣言して下さい。
\usepackage{wrapfig}
%\usepackage{amssymb}
%\usepackage{mb}
%\DeclareGraphicsRule{.tif}{png}{.png}{`convert #1 `dirname #1`/`basename #1 .tif`.png}
\usepackage{lineno}

% ===== my personal definitions ==================================
% ここに、自分のよく使う記号などを定義して下さい。
\newcommand{\klpionn}{K_L \to \pi^0 \nu \overline{\nu}}
\newcommand{\kppipnn}{K^+ \to \pi^+ \nu \overline{\nu}}

% ----- 業績リスト用 -------------
\newcommand{\paper}[6]{%
	% paper{title}{authors}{journal}{vol}{pages}{year}
	\item ``#1'', #2, #3 {\bf #4}, #5 (#6).			% お好みに合わせて変えてください。
}

\newcommand{\etal}{\textit{et al.\ }}
\newcommand{\ca}[1]{*#1}	% corresponding author;   \ca{\yukawa}  みたいにして使う
\newcommand{\invitedtalk}{招待講演}

\newcommand{\yukawa}{H.~Yukawa}					% no underline
%\newcommand{\yukawa}{\underline{\underline{H.~Yukawa}}}	% with 2 underlines
\newcommand{\tomonaga}{S.~Tomonaga}

\newcommand{\prl}{Phys.\ Rev.\ Lett.\ }		% よく使う雑誌も定義すると楽

% ===== 欄外メモ ==================
\newcommand{\memo}[1]{\marginpar{#1}}
%\renewcommand{\memo}[1]{}	% 全てのメモを表示させないようにするには、行頭の"%"を消す

%\input{../../sample/simple/contents}	% skip
\input{pieces/hook5} % pieces

\begin{document}
\input{pieces/hook7} % pieces
%#Split: 01_purpose_plan  
%#PieceName: p01_purpose_plan
\input{pieces/p01_purpose_plan_00}
\section{研究目的、研究方法など}
%　　　　＜＜最大　1ページ＞＞

%s02_purpose_method_cpd
%begin 研究目的１頁 ＝＝＝＝＝＝＝＝＝＝＝＝＝＝＝＝＝＝＝＝
	 象の卵の研究目的は．．．

	\vspace{1cm}
	\begin{thebibliography}{99}
		\bibitem{teramura} 寺村輝夫、「ぼくは王様 - ぞうのたまごのたまごやき」.
	\end{thebibliography}
%end 研究目的１頁 ＝＝＝＝＝＝＝＝＝＝＝＝＝＝＝＝＝＝＝＝
\input{pieces/p01_purpose_plan_01}

%#Split: 02_why_abroad  
%#PieceName: p02_why_abroad
\input{pieces/p02_why_abroad_00}
\section{国際的な環境で研究することの意義・必要性}
%　　　　＜＜最大　1ページ＞＞

% s07_why_abroad
%begin 外国で研究する事の意義 ＝＝＝＝＝＝＝＝＝＝＝＝＝＝＝＝＝＝＝＝
	外国に行きたいから．．．
%end 外国で研究する事の意義 ＝＝＝＝＝＝＝＝＝＝＝＝＝＝＝＝＝＝＝＝

\input{pieces/p02_why_abroad_01}

%#Split: 03_plan  
%#PieceName: p03_plan
\input{pieces/p03_plan_00}
\section{研究計画}
%　　　　＜＜最大　2ページ＞＞

%s08_plan
\noindent
\mbox{（主要渡航期間前）}

%begin 行く前の計画 ＝＝＝＝＝＝＝＝＝＝＝＝＝＝＝＝＝＝＝＝
	うきうきと旅支度．．．
%end 行く前の計画 ＝＝＝＝＝＝＝＝＝＝＝＝＝＝＝＝＝＝＝＝

\noindent
\mbox{（主要渡航期間中）}

%begin 行っている間の計画 ＝＝＝＝＝＝＝＝＝＝＝＝＝＝＝＝＝＝＝＝
	旅の恥はかき捨て、何でも見てやろう．．．
%end 行っている間の計画 ＝＝＝＝＝＝＝＝＝＝＝＝＝＝＝＝＝＝＝＝

\noindent
\mbox{（主要渡航期間後）}

%begin 帰ってからの計画 ＝＝＝＝＝＝＝＝＝＝＝＝＝＝＝＝＝＝＝＝
	思い出のアルバム作り．．．
%end 帰ってからの計画 ＝＝＝＝＝＝＝＝＝＝＝＝＝＝＝＝＝＝＝＝
%========================= 
\input{pieces/p03_plan_01}

%#Split: 99_tail
\input{pieces/hook9} % pieces
\end{document}

